




%PTA introduction
Low frequency gravitational wave (GW) detection is one of the core goals of the global pulsar timing array (PTA) effort. 
PTA experiments are decade-long campaigns monitoring the time of arrival of an array of milli-second pulsars, searching for de
A prime target for PTA GW searches is the nanohertz stochatic GW background (SGWB) produced by the cosmological population of merging massive black hole binaries (MBHB). 

% about MeerKAT
The MeerKAT radio telescope is a $64$ dish interferometer located in the Karoo in the Northern Cape of South Africa.
MeerTime is the MeerKAT pulsar timing program~\citep{MeerTimeBailesEtAl:2020,MeerTimeBailesEtAl:2016}. 
MeerTime pulsar science includes study of relativistic pulsar binaries~\citep{MeerTimeRelBin:2021},
pulsar discovery in globular clusters~\citep{MeerTimeGlobularClusterPSRs:2021}, study of pulsar phenomena in the $1000$ pulsar array~\citep{}, and the MeerTime PTA, which is the focus of this work. 


%PTA current status
PTA searches for the nano-hertz stochastic gravitational wave background are now reaching astrophysically interesting sensitivities. 
Over the years of observations several upper limits on the SGWB amplitude have been placed by the Parkes PTA~\citep[PPTA][]{}, European PTA~\citep[EPTA][]{}, North American Gravitational-wave Observatory~\citep[NANOGrav][]{}. 
Together these three PTAs and the India PTA form the International PTA~\citep[IPTA][]{}, combining datasets to improve SGWB search sensitivity~\cite{}. 
Recent datasets are now showing evidence for a common noise process found by the PPTA, NANOGrav..... as well as the IPTA. 
Future observations will be used to confirm whether or not this signal the signature of a nanohertz SGWB. 
If the signal is indeed real, it will provide insights into the merging population of MBHBs. 

The MeerTime program enables observations of a large number 
The MeerKAT facility provides 

New facilities like MeerTime provide [,,,,, great thigns for pta sensitivity].







In this work we estimate the sensitivity of the MeerTime PTA and test whether small adjustments to the observation strategy can provide sensitivity improvements. 
In Section~\ref{sec:snr} we describe how the PTA S/R is computed.
In Section~\ref{sec:data} we describe the MeerTime data used in this work and the anticipated sensitivity for the current observation strategy. 
We investigate small adjustments to the observing schedule in Section~\ref{sec:shuffle} and present results in Section~\ref{sec:shuffleResults}.
Finally, in Section we see what can be gained from an additional 20\% allocation in observing time?



